% ─────────────────────────────────────────────────────────────────────────────
% FiberRing Routing Evaluation — Paper Draft v0.1
% Authors: Manuel Israel Cázares (Bytepro AI) · Wenlin Zhang (NUS / Omega Institute)
% Title: subject to Wenlin's approval — see comment below
% ─────────────────────────────────────────────────────────────────────────────

\documentclass[11pt]{article}
\usepackage[margin=1in]{geometry}
\usepackage{amsmath,amssymb}
\usepackage{booktabs}
\usepackage{hyperref}
\usepackage{microtype}
\usepackage{listings}
\usepackage{xcolor}
\usepackage{graphicx}
\usepackage[T1]{fontenc}
\usepackage[numbers,sort&compress]{natbib}

% ── Title block ──────────────────────────────────────────────────────────────
% PENDING WENLIN APPROVAL: title, affiliations
\title{Mechanism-level routing failure in LLMs\\
       over Lean-verified algebraic structures}

\author{
  Manuel Israel C\'azares\textsuperscript{1} \and
  Wenlin Zhang\textsuperscript{2}
}

\date{
  \textsuperscript{1}Bytepro AI, Mazatl\'an, M\'exico\\
  \textsuperscript{2}National University of Singapore / Omega Institute\\[6pt]
  \small\textit{Correspondence:} hello@bytepro.ai
}

\begin{document}
\maketitle

% ─────────────────────────────────────────────────────────────────────────────
\begin{abstract}
We present an empirical study of structural routing failure in large language
models (LLMs) over a formally verified algebraic corpus. The task requires
selecting the correct proof-mechanism label---from a fixed closed set of 13
templates---for compact mathematical objects drawn from the FiberRing
formalization in Lean~4~\cite{omega2026automath}, where each evaluation item
is anchored to a Lean-verified artifact and assigned a proof-mechanism label
from the corresponding named certificate family.

Our central finding is a \textbf{mechanism-level routing ceiling}: under
blind conditions (no structural cue), gpt-oss-120b achieves $80.3\%$
template accuracy on 22 clean-anchor FiberRing items
($n=66$ evaluations; temperature~$=0$, seed~$=0$, reasoning~$=\mathrm{low}$),
while Llama~3.3~70B reaches $68.2\%$.
Exposing the Lean-verified verdict as a cue (Condition~A2) raises accuracy
to $90.9\%$ and $81.8\%$ respectively~---~a gap of $+10.6$ and $+13.6$
percentage points that we term \textbf{cue-induced routing uplift}.

The dominant failure mode is a CRT\,$\to$\,ring-equivalence misroute:
gpt-oss-120b misroutes 7 of 12 CRT items ($58.3\%$) in the blind condition,
and zero in the cue-conditioned condition. We show this error has a direct
correlate in the Lean corpus: every CRT row is presented as a ring
equivalence at the outer type level, while the finer mechanism is the coprime
Fibonacci-modulus factorization followed by \texttt{ZMod.chineseRemainder}.
The model identifies the broad
algebraic type but fails to resolve the finer proof mechanism without the
cue.

Two \textbf{cue-resistant} failure modes persist in both conditions
across both models:
\begin{itemize}
  \item \texttt{stable{-}value{-}map} $\to$ \texttt{ring{-}structure}
  \item \texttt{prime{-}power{-}fiber{-}specialization} $\to$
        \texttt{ring{-}equivalence}
\end{itemize}
These errors identify two different boundary phenomena in the Lean
formalization. The \texttt{stable{-}value{-}map}$\,{\to}\,$\texttt{ring{-}structure}
error is a clean formal distinction with strong surface-language overlap:
map-preservation rows mention addition, multiplication, zero, one, and ring
homomorphism, but their proof role is transport through \texttt{toZMod}, not
construction of the ring instance on $X_m$. The
\texttt{prime{-}power{-}fiber{-}specialization}$\,{\to}\,$\texttt{ring{-}equivalence}
error is a true granularity boundary: $X_4 \cong \mathbb{Z}/8\mathbb{Z}$ is
broadly a ring equivalence, but the intended mechanism label is the direct
prime-power specialization of the general stable-value equivalence, with no
nontrivial coprime CRT split.

A cross-model dissociation in Llama~3.3~70B is notable: verdict accuracy
is identical in both conditions ($95.5\%$), while template accuracy
improves $13.6$~pp with the cue. This confirms that truth inference and
proof-mechanism classification are separable capacities---consistent with
the Lean formalization, where a statement's truth value and its certificate
family are encoded at different layers.

These findings extend the \textbf{router hypothesis} introduced in
C\'azares~(2026)~\cite{sair2026} from equational reasoning to formal
algebraic structures: LLMs act as structural classifiers routing over
broad mathematical types, and fail specifically at the level of
fine-grained proof mechanisms. The full evaluation pipeline, prompt pack,
and results are available in the project repository at
\url{https://github.com/bytepro-ai/fiber-routing-eval}.
\end{abstract}

% ─────────────────────────────────────────────────────────────────────────────
\section{Introduction}
\label{sec:intro}
% DRAFT PLACEHOLDER — to be written after Wenlin approves abstract

\section{Related Work}
\label{sec:related}
% DRAFT PLACEHOLDER

\section{Dataset}
\label{sec:dataset}
% DRAFT PLACEHOLDER — FiberRing corpus, Set A design, 22-item schema

\section{Methodology}
\label{sec:methodology}
% DRAFT PLACEHOLDER — A1/A2 conditions, label set, scoring protocol

\section{Results}
\label{sec:results}
% DRAFT PLACEHOLDER — headline numbers, confusion matrix

\section{Analysis}
\label{sec:analysis}
% Wenlin's Section 6 mathematical interpretation goes here

\section{Discussion}
\label{sec:discussion}
% DRAFT PLACEHOLDER — router hypothesis extension, limitations

\section{Conclusion}
\label{sec:conclusion}
% DRAFT PLACEHOLDER

% ─────────────────────────────────────────────────────────────────────────────
\bibliographystyle{plainnat}
\bibliography{references}

\end{document}